\section{Introduction}
\label{sec:introduction}

Capacitors are reliability-critical components of power electronic converters, and their condition strongly affects converter availability \cite{wang2014reliability}. However, capacitors degrade under sustained electrical and thermal stress, and the resulting capacitance loss and equivalent series resistance (ESR) growth in turn alter electrical stress, thermal loss, and filtering performance; both quantities are established degradation precursors for electrolytic and film capacitors \cite{wang2014reliability,yao2024film}. Since neither parameter is directly measurable during operation, many online monitoring methods have been developed, spanning ripple analysis, injected excitation, circuit reconfiguration, observers, recursive identification, and data-driven processing \cite{dang2020review,nathan2023review,zhao2021overview}. The suitability of a given scheme depends on the capacitor's circuit role and on which electrical variables are already measured.

The energy-transfer capacitor in a \Cuk{} converter is distinct from a dc-link or output-filter capacitor: it lies in the main power-transfer path and alternately carries the full inductor currents $+i_{L1}$ and $-i_{L2}$. Every switching transition therefore commutates the capacitor current between two signed values, and normal operation supplies this bidirectional excitation at the switching frequency without any diagnostic perturbation. Treating this component as a generic two-terminal capacitor sampled asynchronously discards this topology-specific information.

Existing joint C--ESR identification methods approach the problem from other circuit roles. Hasegawa \emph{et al.} \cite{hasegawa2018dclink} estimate dc-link capacitance and ESR by separating ripple-frequency components in an inverter dc link; however, the separation is tied to the ripple bands of that topology. Yao \emph{et al.} \cite{yao2015buck} identify ESR and capacitance of a buck output capacitor from the pulsewidth-modulation (PWM) signal and output voltage without a current sensor, and related sensorless formulations extend this route \cite{tang2019sensorless}; however, the underlying within-cycle relation presumes the unidirectional filter current of the buck output stage and does not carry over to a capacitor whose current reverses in every subinterval. Observer-based schemes \cite{barbulescu2022observer,barbulescu2023implementation,caiman2024quasionline} obtain joint estimates through a switched resistive network; however, they require an added discharge network or restrict monitoring to quasi-online intervals. Ribeiro \emph{et al.} \cite{ribeiro2025inherent,ribeiro2024injection,ribeiro2023wavelet} extract inherent or injected spectral features for neutral-point-clamped (NPC) dc-link capacitors; however, injection adds loss, and the inherent-signal mapping is specific to the dc-link spectrum. Recent extensions include sensorless boost ESR estimation, optimized injection for buck capacitors, and large-signal C--ESR estimation \cite{leong2025boostesr,zhang2025injection,zhang2025largesignal}. Experimentally validated dc-link monitoring continues to advance through maximum-likelihood traction-inverter estimation, adaptive-forgetting recursive least squares (RLS) for cascaded-H-bridge converters, discharge-profile and grey-box transient methods, multilevel-converter ESR--C schemes, and DSP-efficient spectral precursors \cite{wu2025mlrail,wang2024chb,wei2024discharge,zhao2023greybox,deng2020mmc,asoodar2024mmc,sundararajan2020goertzel}, and wavelet/Kalman and RLS/Kalman formulations provide further estimation routes \cite{amaral2024waveletkf,zhang2025haar,reichbach2016rls,chen2018vffrls}. For the \Cuk{} converter itself, fractional-order modeling, converter fault diagnosis, sensor-fault detection, and prognostics have been studied \cite{wang2021fractionalcuk,prathiba2023cuk,ajra2024sensorfault,samie2015prognostic}; however, none of these works exploits the transfer-capacitor current commutation for joint C--ESR identification.

The noninvasive joint schemes of \cite{hasegawa2018dclink,yao2015buck,ribeiro2025inherent,amaral2024waveletkf} and the observer route of \cite{barbulescu2022observer} are the closest related methods; the comparison remains categorical rather than numerical because the converter topologies, sensing sets, and reporting protocols differ across studies. Relative to all of them, this work is distinguished by exploiting the $+i_{L1}\leftrightarrow-i_{L2}$ commutation of the energy-transfer capacitor itself, with the terminal voltage, the two inductor currents, and the switching state and timestamp as the only measurements. The residual gap is structural rather than algorithmic: no reported method converts the \Cuk{} transfer mechanism itself into implementable identification information.

Therefore, in this paper, topology-synchronous decoupled observations are proposed that perform exactly this conversion: local fits to timestamped samples in guarded pre- and post-edge windows reconstruct the ESR-dominant terminal-voltage discontinuity at the switching instant, and an ESR-compensated charge balance inside one switching state yields a capacitance-dominant row. In the common-kernel factorial of Section~\ref{sec:results}, the proposed observations alone reduce mean capacitance error from $7.0366\%$ to $0.3727\%$ with the RLS kernel unchanged. On these observations, a topology-synchronous supervised-reset Kalman estimator (\TSSRKE) is constructed, which restores gain after abrupt health changes while leaving steady-state behavior unchanged. The design is evaluated on 48 blind cases across four noise profiles and four timing skews, with abrupt-step, degradation-ramp, closed-loop-regulation, light-load, and TMS320F28379D device-realistic scenarios.

The main contributions of this paper are as follows.
\begin{enumerate}
\item Topology-synchronous decoupled observations that interpret the normal \Cuk{} switching sequence as topology-induced health excitation and convert it, through timestamp-reconstructed edge fits and safe-window charge balances, into ESR-dominant and capacitance-dominant rows.
\item Joint finite-window identifiability established from the raw observation pair, with positive per-direction information bounds for the implemented sequential stream, under valid continuous-conduction-mode conditions.
\item An observation--estimator factorial inside a predeclared blind verification protocol (48 cases, a conditional Cram\'er--Rao benchmark and sensitivity analysis, and a TMS320F28379D device-realistic simulation) demonstrating that the observation design, rather than added estimator complexity, carries the dominant accuracy gain.
\item A supervised-reset extension (\TSSRKE) that restores gain after abrupt health changes through a predeclared two-time-scale supervisor over the accepted \emph{and} gate-rejected normalized innovations; its components are classical, and the claimed novelty is confined to their direction-specific construction on the gated topology-synchronous rows (Section~\ref{sec:srke}).
\end{enumerate}

The rest of this paper is organized as follows. Section~\ref{sec:modeling} develops the switched model of the energy-transfer capacitor and its health signatures. Section~\ref{sec:observations} presents the topology-synchronous decoupled observation framework. Finite-window identifiability and information bounds are established in Section~\ref{sec:identifiability}. Section~\ref{sec:estimators} presents the estimation kernels and the supervised-reset realization. Section~\ref{sec:protocol} describes the verification methodology. In Section~\ref{sec:results}, simulation results are reported. Section~\ref{sec:device} assesses device-realistic feasibility on the TMS320F28379D. Section~\ref{sec:discussion} discusses scope and limitations. Finally, Section~\ref{sec:conclusion} concludes the paper.
